% Don't touch this %%%%%%%%%%%%%%%%%%%%%%%%%%%%%%%%%%%%%%%%%%%
\documentclass[11pt]{article}
\usepackage{fullpage}
\usepackage[left=1in,top=1in,right=1in,bottom=1in,headheight=3ex,headsep=3ex]{geometry}
\usepackage{graphicx}
\usepackage{float}
\usepackage{quoting}

\setlength{\parindent}{0em}
\setlength{\parskip}{1em}

\newcommand{\blankline}{\quad\pagebreak[2]}
%%%%%%%%%%%%%%%%%%%%%%%%%%%%%%%%%%%%%%%%%%%%%%%%%%%%%%%%%%%%%%

% Modify Course title, instructor name, semester here %%%%%%%%

\title{IDS-101-10: College Colloquium: Thinking Machines}
\author{Calvin Deutschbein (they) and Addie Gage (she)}
\date{Willamette University, Fall 2025}

%%%%%%%%%%%%%%%%%%%%%%%%%%%%%%%%%%%%%%%%%%%%%%%%%%%%%%%%%%%%%%

% Don't touch this %%%%%%%%%%%%%%%%%%%%%%%%%%%%%%%%%%%%%%%%%%%
\usepackage[sc]{mathpazo}
\linespread{1.05} % Palatino needs more leading (space between lines)
\usepackage[T1]{fontenc}
\usepackage[mmddyyyy]{datetime}% http://ctan.org/pkg/datetime
\usepackage{advdate}% http://ctan.org/pkg/advdate
\newdateformat{syldate}{\twodigit{\THEMONTH}/\twodigit{\THEDAY}}
\newsavebox{\MONDAY}\savebox{\MONDAY}{Mon}% Mon
\newcommand{\week}[1]{%
%  \cleardate{mydate}% Clear date
% \newdate{mydate}{\the\day}{\the\month}{\the\year}% Store date
  \paragraph*{\kern-2ex\quad #1, \syldate{\today} - \AdvanceDate[4]\syldate{\today}:}% Set heading  \quad #1
%  \setbox1=\hbox{\shortdayofweekname{\getdateday{mydate}}{\getdatemonth{mydate}}{\getdateyear{mydate}}}%
  \ifdim\wd1=\wd\MONDAY
    \AdvanceDate[7]
  \else
    \AdvanceDate[7]
  \fi%
}
\usepackage{setspace}
\usepackage{multicol}
%\usepackage{indentfirst}
\usepackage{fancyhdr,lastpage}
\usepackage{url}
\pagestyle{fancy}
\usepackage{hyperref}
\usepackage{lastpage}
\usepackage{amsmath}
\usepackage{layout}

\lhead{}
\chead{}
%%%%%%%%%%%%%%%%%%%%%%%%%%%%%%%%%%%%%%%%%%%%%%%%%%%%%%%%%%%%%%

% Modify header here %%%%%%%%%%%%%%%%%%%%%%%%%%%%%%%%%%%%%%%%%
\rhead{\footnotesize Thinking Machines}

%%%%%%%%%%%%%%%%%%%%%%%%%%%%%%%%%%%%%%%%%%%%%%%%%%%%%%%%%%%%%%
% Don't touch this %%%%%%%%%%%%%%%%%%%%%%%%%%%%%%%%%%%%%%%%%%%
\lfoot{}
\cfoot{\small \thepage/\pageref*{LastPage}}
\rfoot{}

\usepackage{array, xcolor}
\usepackage{color,hyperref}
\definecolor{clemsonorange}{HTML}{EA6A20}
\hypersetup{colorlinks,breaklinks,linkcolor=clemsonorange,urlcolor=clemsonorange,anchorcolor=clemsonorange,citecolor=black}

\begin{document}

\maketitle

\blankline

\begin{tabular*}{.93\textwidth}{@{\extracolsep{\fill}}lr}

%%%%%%%%%%%%%%%%%%%%%%%%%%%%%%%%%%%%%%%%%%%%%%%%%%%%%%%%%%%%%%

% Modify information %%%%%%%%%%%%%%%%%%%%%%%%%%%%%%%%%%%%%%%%%
E-mail: \texttt{ckdeutschbein@willamette.edu} & Web: \href{https://cd-public.github.io/courses/soc}{\tt\bf cd-public.github.io/}  \\

 Office Hours: MWF 11:30-12:00 \& Discord TTh  &  Lecture: MWF 12:00-1:00 \\

 Office: Ford 307 & Chase Hour: Th 12:00-1:00  \\
 & \\
\hline
\end{tabular*}

\vspace{5 mm}

% First Section %%%%%%%%%%%%%%%%%%%%%%%%%%%%%%%%%%%%%%%%%%%%

\section*{Course Description}

\subsection*{College Colloquium}

Topical seminars designed to pursue significant issues and questions of special interest to instructors and students. Seminars invite students into the intellectual life of the university, model rigorous engagement, and help them develop qualities of good scholarship -- effective writing, careful reading, critical thinking, and cogent argumentation. Seminars do not count toward majors or minors. Required for all entering firstyear students.

\subsection*{Thinking Machines}

Artifical intelligence, machine learning, big data, and large scale computing have captured the collective imaginations of millions people across a range of technical expertise. Yet each of these buzzwords is rooted in a central promise: that machines may think.

Guided by student inquiry, we will learn how to construct thinking - or the appearance thereof - from basic computation. Using the shared language of computer code, we use programming as a way to bridge the gap between thinking humans and thinking machines. With this shared experience, we will be able to reflect more thoughtfully on the implications of thinking machines on the self, culture, and society. 

\section*{About Us}
Calvin Deutschbein is a 
fifth-year professor of computer science and colloquia instructor.
%They moved to Oregon from North Carolina in 2021 after being involved in a number of worker campaigns in the US South, including faculty unionization at Elon University and a number of moderate policing reforms in Chapel Hill.

Addie Gage is a fourth-year student at Willamette studying Computer Science and Data Science in the MS Data Science program, and a STEM fellow.

% Second Section %%%%%%%%%%%%%%%%%%%%%%%%%%%%%%%%%%%%%%%%%%%

\section*{Required Materials}

Required materials for a given class will be available on the \href{https://cd-public.github.io/scicom}{course webpage}. All course materials will be made available at no cost to the student.


% Third Section %%%%%%%%%%%%%%%%%%%%%%%%%%%%%%%%%%%%%%%%%%%


\section*{Accessability}

I will make every effort to ensure all coursework and materials are accessible to all students, including working with on-campus specialists. However, there is always room for improvement. I always appreciate hearing from students about how I can make the course more accessible, so please reach out if there is something I can be doing better!

% Fourth Section %%%%%%%%%%%%%%%%%%%%%%%%%%%%%%%%%%%%%%%%%%%

\section*{Course Objectives}
Per the college, we learn and practice techniques to:
\begin{itemize}
\item Critically examine information and/or texts (written, oral, artistic, or quantitative) by identifying central ideas or arguments, making inferences, questioning underlying assumptions, and assessing evidence.
\item Learn to contribute to a constructive classroom climate through participation in thoughtful, informed, and responsive discussion, effective speaking, active listening, and the development of an iterative group process of critical analysis and interpretation.
\item Effectively formulate ideas and arguments, develop them through an iterative process, and express them clearly and persuasively via linguistic, artistic, and/or quantitative modes of communication.
\end{itemize}
Additionally, I will attempt to construct the class in a manner that these techniques specifically create an inclusive academic environment.

% Fifth Section %%%%%%%%%%%%%%%%%%%%%%%%%%%%%%%%%%%%%%%%%%%

\section*{Course Structure}

The course will be composed of lecture, demonstration, and labs on scientific computing, and out-of-class writing of both technical and non-technical nature.

Students will also read and respond to each other's writing within class.

\subsection*{Class Structure}

Classes are scheduled for Monday, Wednesday, and Friday at 12:00 noon. There is a special supporting CHASE session on Thursdays at 12:00 noon. The lecture schedule is on the \href{https://cd-public.github.io/scicom}{course webpage}.

\subsubsection*{Papers/Projects}

There will be both a midterm paper/project and final paper/project for this course. Though we term them papers, they can more properly be considered scientific publications. The midterm will be a source code repository on GitHub with both Python code and Markdown documentation. The final will be a GitHub Pages website created using the Quarto platform for scientific publication. These tools will be taught during the course.

The midterm will be due 10 October at 12:00 PM. 
The final paper will be due 14 November at 12:00 PM.


\subsection*{Feedback and Grading}

\subsubsection*{Exercises}

\begin{itemize}
    \item For an A (100\%)
    \begin{itemize}
        \item Code that meets requirements.
        \item On time.
    \end{itemize}
    \item For a B (90\%)
    \begin{itemize}
        \item Code that runs without error.
        \item On time.
    \end{itemize}
    \item Erroneous code is an F (60\%).
    \item Late work is a zero (0\%).
\end{itemize}

\subsubsection*{Papers}

Graded on \href{https://apcentral.collegeboard.org/media/pdf/ap-english-language-and-composition-frqs-1-2-3-scoring-rubrics.pdf}{AP Scale} as an argument essay in support (or against) scientific computing.  If I take take more than one week to return grades, I provide ``grading amnesty" of up to a letter grade.

\subsubsection*{Weighting}

There will be 9 exercises of equal weight, and midterm weighted as 3 exercises and a final weighted as 6 exercises. I will drop the three lowest (which could be up to half the final or the entire midterm) and take the average of the top 15 scores. 

Students should expect to earn to perfect scores on exercises unless skipping class, which essentially amounts to an attendance incentive of 60\% of the final grade.

\input{college}

\end{document}

